% !PRINTED 179
\weberpage{179}

Thus, for example, every linear function
\[
y = \overset{i}{\Sigma} a_i x_i
\]
passes into a function of the same kind
\[
y = \overset{i}{\Sigma} a_i' x_i'
\]
in which
\weq{3}{a_k' = \overset{i}{\Sigma} a_i \alpha_i^{(k)}.}
If we consider $m$ such linear functions
\[
y_\nu = \overset{i}{\Sigma} a_{i,\nu} x_i \qquad \nu = 1, 2 \ldots m,
\]
we obtain $m$ transformed functions
\[
y_\nu = \overset{i}{\Sigma} a_{i,\nu}' x_i' \qquad \nu = 1, 2 \ldots m.
\]
The determinants of the systems $y_\nu, y_\nu'$
\[
\Delta = \Sigma \pm a_{1,1} a_{2,2} \ldots a_{m,m}, \qquad \Delta' = \Sigma \pm a_{1,1}' a_{2,2}' \ldots a_{m,m}'
\]
stand, by the multiplication law of determinants, in the relation
\weq{4}{\Delta' = r \Delta,}
so that the one does not vanish without the other.

\begin{center}
\S.\,56.

\gesperrt{Quadratic Forms.}
\end{center}

Of particular importance is the transformation of the homogeneous functions of the second degree, or the quadratic forms. We denote them, as already in the first section, thus:
\weq{1}{\varphi(x_1, x_2 \ldots x_m) = \overset{i,k}{\Sigma} a_{i,k} x_i x_k, \qquad a_{i,k} = a_{k,i}.}
By the substitution \S.\,55, (1), $\varphi$ passes into
\weq{2}{\Phi(x_1', x_2' \ldots x_m') = \overset{i,k}{\Sigma} a_{i,k}' x_i' x_k'.}
If we put $x + \xi$ for $x$ and likewise $x' + \xi'$ for $x'$, so that the $\xi'$ stand in the same relation to the $\xi$ as the $x'$ to the $x$, then, if for abbreviation we set
\[
\tfrac{1}{2}\varphi'(x_i) = u_i, \qquad \tfrac{1}{2}\Phi'(x_i') = u_i'
\]
we obtain, by comparison of the terms of the first dimension (\S.\,16):
%CONT>>
